\documentclass[11pt]{amsart}

\usepackage[T1]{fontenc}
\usepackage{lmodern}
\usepackage{microtype}
\usepackage{mathtools,amssymb,amsthm}
\usepackage{array}
\usepackage[hidelinks]{hyperref}

\emergencystretch=2em

\hypersetup{
  pdftitle={Two Zonotopes with the Same Arithmetic Matroid and Different Graded Ehrhart Series}
}

\newtheorem{theorem}{Theorem}
\newtheorem{lemma}[theorem]{Lemma}
\theoremstyle{definition}
\newtheorem*{cpq}{Question \textup{(Crowley--Partida \cite{CP})}}

\newcommand{\Z}{\mathbb Z}
\newcommand{\R}{\mathbb R}
\newcommand{\C}{\mathbb C}
\newcommand{\Orb}{\operatorname{Orb}}
\newcommand{\Hilb}{\operatorname{Hilb}}
\newcommand{\grI}{\operatorname{gr}}
\newcommand{\Zone}{Z_1}
\newcommand{\Ztwo}{Z_2}

\title[Equal arithmetic matroid, different graded Ehrhart series]
{Two Zonotopes with the Same Arithmetic Matroid\\ and Different Graded Ehrhart Series}

\date{}

\subjclass[2020]{52C07, 05B35, 52B20, 05A15}
\keywords{graded Ehrhart theory, zonotope, arithmetic matroid, arithmetic Tutte polynomial,
harmonics, counterexample}

\begin{document}

\begin{abstract}
Crowley and Partida ask whether the graded Ehrhart theory of a non-unimodular zonotope is an
invariant of its associated arithmetic matroid. The answer is no, by a pair of lattice
parallelograms in $\Z^2$: the zonotopes $\Zone$ generated by $(1,0),(1,5)$ and $\Ztwo$ generated
by $(1,0),(2,5)$ have literally equal arithmetic matroids, hence the same arithmetic Tutte
polynomial $x^2+4$ and the same ungraded lattice point counts $5m^2+2m+1$, yet
\[
 i_{\Zone}(1;q)=1+2q+2q^2+q^3+q^4+q^5,
 \qquad
 i_{\Ztwo}(1;q)=1+2q+3q^2+2q^3 .
\]
The witness is exhibited in full and the separation is proved by hand at $m=1$.
\end{abstract}

\maketitle

\section{The question}

Let $S\subset\Z^d$ be finite, let $I(S)\subseteq\C[x_1,\dots,x_d]$ be its vanishing ideal, and
let $\grI I(S)$ be the homogeneous ideal generated by the top-degree homogeneous components of
all elements of $I(S)$. Put
\[
 \Orb(S)=\C[x_1,\dots,x_d]/\grI I(S),
\]
a graded ring of total dimension $|S|$. For a lattice polytope $P\subseteq\R^d$ and an integer
$m\ge0$, the \emph{graded Ehrhart series} of $P$ is
\begin{equation}\label{eq:def}
 i_P(m;q)=\Hilb\bigl(\Orb(mP\cap\Z^d);q\bigr),
\end{equation}
so that $i_P(m;1)=|mP\cap\Z^d|$ is the ordinary Ehrhart polynomial. This is the notion of Reiner
and Rhoades \cite{RR}; \eqref{eq:def} is Definition~2.4 of Crowley and Partida \cite{CP}.

Let $v_1,\dots,v_n\in\Z^d$ and let $Z=Z(v_1,\dots,v_n)=\{\sum_it_iv_i:0\le t_i\le1\}$ be the
zonotope they generate. Its \emph{arithmetic matroid} is the represented matroid on
$E=\{1,\dots,n\}$ together with the \emph{multiplicity function}
\begin{equation}\label{eq:mult}
 \mu(S)=\bigl[\,\bigl(\operatorname{span}_\R\{v_i:i\in S\}\cap\Z^d\bigr):\textstyle\sum_{i\in S}\Z v_i\,\bigr],
\end{equation}
the index of the lattice generated by $S$ inside the saturation of its span; on a basis this is
$|\det|$. D'Adderio and Moci \cite{DM} proved that the ungraded counts are determined by these
data: $|mZ\cap\Z^d|=m^{d}\,T_M\!\left(\frac{m+1}{m},1\right)$, where $T_M$ is the arithmetic
Tutte polynomial. Crowley and Partida prove that for a \emph{unimodular} zonotope the graded
series is likewise matroidal, and then ask:

\begin{cpq}
Is the graded Ehrhart theory of a non-unimodular zonotope $Z$ an invariant of its associated
arithmetic matroid? If so, is there a generalization of \textup{[}their proposition on Ehrhart
polynomials\textup{]} to arithmetic matroids?
\end{cpq}

\noindent The wording is quoted from the source of \cite{CP}, where the question stands in the
subsection \emph{The graded Ehrhart theory of non-unimodular zonotopes}. Only the first clause is
at issue here. The same subsection states the expectation:
\begin{quote}
``[\,\dots\,] it is still possible that the $q$-lattice point counts of an arbitrary
zonotope are an arithmetic matroid invariant. It would be quite surprising to find two zonotopes
$Z$ and $Z'$ whose arithmetic matroids are the same but $i_Z(m;q)\ne i_{Z'}(m;q)$ for some
integer $m\ge0$.''
\end{quote}

\section{The witness}

\begin{theorem}\label{thm:main}
Let
\[
 A_1=\begin{pmatrix}1&1\\0&5\end{pmatrix},
 \qquad
 A_2=\begin{pmatrix}1&2\\0&5\end{pmatrix},
\]
and let $\Zone=Z\bigl((1,0),(1,5)\bigr)$ and $\Ztwo=Z\bigl((1,0),(2,5)\bigr)$ be the zonotopes
generated by their columns. Then:
\begin{enumerate}
\item[(a)] the arithmetic matroids of $\Zone$ and $\Ztwo$ on the ground set $E=\{1,2\}$ are
\emph{equal}, not merely isomorphic: both are $U_{2,2}$ with
$\mu(\varnothing)=\mu(\{1\})=\mu(\{2\})=1$ and $\mu(\{1,2\})=5$, hence both have arithmetic
Tutte polynomial $T_M(x,y)=x^2+4$ and, for every $m\ge0$,
$|m\Zone\cap\Z^2|=|m\Ztwo\cap\Z^2|=5m^2+2m+1$;
\item[(b)]
\[
 i_{\Zone}(1;q)=1+2q+2q^2+q^3+q^4+q^5
 \ne
 1+2q+3q^2+2q^3=i_{\Ztwo}(1;q).
\]
\end{enumerate}
Consequently the answer to the first clause of the question above is \textbf{no}.
\end{theorem}

The two lattice point sets at $m=1$ are, in full,
\begin{align}
 S_1&=\Zone\cap\Z^2=\{(0,0),(1,0),(1,1),(1,2),(1,3),(1,4),(1,5),(2,5)\},\label{eq:S1}\\
 S_2&=\Ztwo\cap\Z^2=\{(0,0),(1,0),(1,1),(1,2),(2,3),(2,4),(2,5),(3,5)\}.\label{eq:S2}
\end{align}
Both have eight elements. Six of the eight points of $S_1$ are collinear; the largest collinear
subset of $S_2$ has three points. That difference in configuration is the whole phenomenon, and
Section~\ref{sec:proof} turns it into a proof.

Part~(a) is immediate: all four generating columns are primitive, so
$\mu(\{1\})=\mu(\{2\})=1$ for both, and $|\det A_1|=|\det A_2|=5$, so $\mu(\{1,2\})=5$ for both.
Hence
$T_M(x,y)=(x-1)^2+2(x-1)+5=x^2+4$, and the D'Adderio--Moci identity gives
$m^2\bigl(((m+1)/m)^2+4\bigr)=5m^2+2m+1$ for both. (Equivalently, by Pick's theorem: area $5m^2$,
boundary points $B=4m$ because all four edges are primitive, interior points $I=5m^2-2m+1$.)

\section{Proof of Theorem~\ref{thm:main}(b)}\label{sec:proof}

For a finite $S\subset\Z^2$ and $j\ge0$ let $r_j(S)$ be the rank of the evaluation matrix of the
monomials of degree at most $j$ on $S$, and let $N_j=\binom{j+2}{2}$ be the number of those
monomials. Write $W_j$ for the space of degree-$j$ homogeneous top components of elements of
$I(S)$ of degree at most $j$.

\begin{lemma}\label{lem:diff}
$\dim\Orb(S)_j=r_j(S)-r_{j-1}(S)$ for all $j\ge0$, with $r_{-1}=0$.
\end{lemma}

\begin{proof}
Taking the degree-$j$ component is a linear surjection from $\{f\in I(S):\deg f\le j\}$ onto
$W_j$ with kernel $\{f\in I(S):\deg f\le j-1\}$, so
$\dim W_j=(N_j-r_j)-(N_{j-1}-r_{j-1})=(j+1)-(r_j-r_{j-1})$. Moreover $xW_{j-1}+yW_{j-1}\subseteq W_j$: if
$f\in I(S)$ has degree $\le j-1$ then $xf\in I(S)$ has degree $\le j$ and its degree-$j$
component is $x$ times the degree-$(j-1)$ component of $f$. Hence the homogeneous ideal generated
by $\bigcup_iW_i$ has degree-$j$ part exactly $W_j$, that is
$\bigl(\grI I(S)\bigr)_j=W_j$, whence
$\dim\Orb(S)_j=\dim\C[x,y]_j-\dim W_j=(j+1)-\dim W_j=r_j-r_{j-1}$.
\end{proof}

\subsection*{$S_1$: a forced line component}
$S_1$ from \eqref{eq:S1} has six points on the line $x=1$, namely $(1,y)$ for $0\le y\le5$, and
the two extra points $(0,0)$ and $(2,5)$. Let $2\le j\le5$ and let $f$ have degree at most $j$
and vanish on $S_1$. Restricting $f$ to $x=1$ gives a polynomial in $y$ of degree at most
$j\le5$ with six distinct roots, hence identically zero; therefore $(x-1)\mid f$. Since $x-1$ is
nonzero at $(0,0)$ and at $(2,5)$, the cofactor $g=f/(x-1)$ has degree at most $j-1$ and
vanishes at those two points, and the conditions $g(0,0)=g(2,5)=0$ are independent on
$\C[x,y]_{\le j-1}$ for $j\ge2$ because $1$ and $x$ separate the two points. Conversely every
such product vanishes on $S_1$. Hence
\begin{align*}
 \dim\{f\in I(S_1):\deg f\le j\}&=N_{j-1}-2,\\
 r_j(S_1)&=N_j-N_{j-1}+2=j+3
 \qquad(2\le j\le5).
\end{align*}
With $r_0=1$ and $r_1=3$ the profile is $r_\bullet=(1,3,5,6,7,8)$, and Lemma~\ref{lem:diff}
gives $\dim\Orb(S_1)_\bullet=(1,2,2,1,1,1)$, i.e.
$i_{\Zone}(1;q)=1+2q+2q^2+q^3+q^4+q^5$. The coefficients sum to $8=|S_1|$.

\subsection*{$S_2$: no conic, and eight separating cubics}
Label the points of \eqref{eq:S2}
\begin{gather*}
 P_1=(0,0),\quad P_2=(1,0),\quad P_3=(1,1),\quad P_4=(1,2),\\
 P_5=(2,3),\quad P_6=(2,4),\quad P_7=(2,5),\quad P_8=(3,5).
\end{gather*}
Clearly $r_0=1$ and $r_1=3$ (the points are not collinear).

\emph{$r_2=6$, that is, no conic passes through $S_2$.} A conic meets a line in at most two
points unless the line is one of its components. The three points $P_2,P_3,P_4$ lie on $x=1$ and
the three points $P_5,P_6,P_7$ lie on $x=2$, so a conic through $S_2$ must be a scalar multiple
of $(x-1)(x-2)$; that polynomial takes the value $(-1)(-2)=2$ at $P_1=(0,0)$, a contradiction.
Since $N_2=6$, the eight points impose $6$ independent conditions on conics.

\emph{$r_3=8$, that is, the eight points impose independent conditions on cubics.} It suffices to
exhibit, for each $i$, a cubic vanishing at the seven points $P_j$, $j\ne i$, and not at $P_i$.
Table~\ref{tab:cubics} does so; each cubic is a product of three lines, and the column ``lines
cover'' shows how the seven points are distributed among them, so the vanishing is visible
without any substitution. The stated value at the omitted point is the product of the three
displayed linear forms evaluated there.

\begin{table}[htbp]
\small
\renewcommand{\arraystretch}{1.2}
\begin{tabular}{@{}clr>{\raggedright\arraybackslash}p{0.40\textwidth}@{}}
\hline
omit & cubic & value & lines cover\\
\hline
$P_1=(0,0)$ & $(x-1)(x-2)(x-3)$ & $-6$ & $x{=}1$: $P_2P_3P_4$; $x{=}2$: $P_5P_6P_7$; $x{=}3$: $P_8$\\
$P_2=(1,0)$ & $(y-2x)(x-2)(y-2x+1)$ & $-2$ & $y{=}2x$: $P_1P_4P_6$; $x{=}2$: $P_5P_7$; $y{=}2x{-}1$: $P_3P_8$\\
$P_3=(1,1)$ & $y\,(x-2)(2y-3x-1)$ & $2$ & $y{=}0$: $P_1P_2$; $x{=}2$: $P_5P_6P_7$; $2y{=}3x{+}1$: $P_4P_8$\\
$P_4=(1,2)$ & $y\,(x-2)(y-2x+1)$ & $-2$ & $y{=}0$: $P_1P_2$; $x{=}2$: $P_5P_6P_7$; $y{=}2x{-}1$: $P_3P_8$\\
$P_5=(2,3)$ & $(y-5)(x-1)(y-2x)$ & $2$ & $y{=}5$: $P_7P_8$; $x{=}1$: $P_2P_3P_4$; $y{=}2x$: $P_1P_6$\\
$P_6=(2,4)$ & $(y-5)(x-1)(3x-2y)$ & $2$ & $y{=}5$: $P_7P_8$; $x{=}1$: $P_2P_3P_4$; $3x{=}2y$: $P_1P_5$\\
$P_7=(2,5)$ & $(y-2x)(x-1)(y-2x+1)$ & $2$ & $y{=}2x$: $P_1P_4P_6$; $x{=}1$: $P_2P_3$; $y{=}2x{-}1$: $P_5P_8$\\
$P_8=(3,5)$ & $x\,(x-1)(x-2)$ & $6$ & $x{=}0$: $P_1$; $x{=}1$: $P_2P_3P_4$; $x{=}2$: $P_5P_6P_7$\\
\hline
\end{tabular}
\medskip
\caption{Eight cubics separating the points of $S_2$, proving $r_3(S_2)=8$.}
\label{tab:cubics}
\end{table}

Hence $r_\bullet(S_2)=(1,3,6,8)$ and $r_3=8=|S_2|$, so by Lemma~\ref{lem:diff}
$\dim\Orb(S_2)_\bullet=(1,2,3,2)$ and $i_{\Ztwo}(1;q)=1+2q+3q^2+2q^3$, again of coefficient sum
$8$. The two series differ already in the coefficient of $q^2$ ($2$ against $3$) and in degree
($5$ against $3$), while part~(a) makes the arithmetic matroids equal. One value of $m$ suffices,
so Theorem~\ref{thm:main} is proved.\hfill\qedsymbol

\medskip

\noindent The graded series are compared here only at $m=1$; the second clause of the question and
the case of three or more generators are not addressed.

\section{Reproduction and verification}\label{sec:verif}

Everything above can be redone by hand from \eqref{eq:S1}, \eqref{eq:S2} and
Section~\ref{sec:proof}; no computation is needed to check Theorem~\ref{thm:main}.

A program \texttt{verify.py} accompanies this note, together with a recorded run
\texttt{verify.output.txt}. It re-derives every numerical claim above from the two generator
matrices alone, in exact integer and rational arithmetic: the multiplicity function by integer
Smith normal form, the lattice point sets by two independent tests cross-checked against
\eqref{eq:S1}, \eqref{eq:S2} and Pick's theorem, and each graded series both by the rank
differencing of Lemma~\ref{lem:diff} and from an explicit basis of $\{f\in I(S):\deg f\le j\}$ and
its degree-$j$ top forms. It confirms every row of Table~\ref{tab:cubics}, including the printed
value of each cubic at its omitted point. All $48$ of its checks pass; several of them concern
material not used above, and its output lists what it does not cover, the bibliographic facts among
them: those are read from the sources, not computed.

\begin{thebibliography}{9}

\bibitem{CP}
S.~Crowley and D.~Partida,
\emph{Graded Ehrhart theory of unimodular zonotopes},
arXiv:2603.07873, 2026.
The question quoted in Section~1 is Question~7.3 there, read from the \LaTeX{} source of the
e-print.

\bibitem{DM}
M.~D'Adderio and L.~Moci,
\emph{Arithmetic matroids, the Tutte polynomial and toric arrangements},
Adv. Math. \textbf{232} (2013), 335--367.

\bibitem{RR}
V.~Reiner and B.~Rhoades,
\emph{Harmonics and graded Ehrhart theory},
\href{https://arxiv.org/abs/2407.06511}{arXiv:2407.06511}, 2024.

\end{thebibliography}

\end{document}
